% Vietnamese translation for OI Public Library
% Source-PDF: IOI中国国家候选队论文/2001/江鹏.pdf
\documentclass[11pt]{article}
\usepackage{oipl-vn}

\title{Từ lời giải một bài toán bàn về xây dựng mô hình và thuật toán luồng mạng}
\author{Jiang Peng}
\date{}

\begin{document}
\maketitle

\origtitle{Network-flow construction and algorithms through one problem}
\sourcepath{IOI China candidate team papers / 2001 / Jiang Peng.pdf}

\section{Dẫn nhập}

\subsection{Nhận thức về thuật toán luồng mạng}

Thuật toán luồng mạng là một loại thuật toán hiệu quả và thực dụng. So với các
thuật toán đồ thị khác, mô hình của nó phức tạp hơn và độ phức tạp lập trình
cũng cao hơn. Nhưng nó tổng hợp được một số thuật toán khác trong đồ thị, ví
dụ đường đi ngắn nhất, nên phạm vi áp dụng rộng hơn. Nhiều khi nó có thể giải
tốt những bài toán mà tìm kiếm hoặc quy hoạch động không xử lý được, và thoạt
nhìn có vẻ giống bài toán NP.

\subsection{Ứng dụng vào bài toán cụ thể}

Khi áp dụng luồng mạng vào một bài toán cụ thể, phần thách thức nhất là xây
dựng mô hình. Việc này không có một khuôn mẫu sẵn để áp vào. Người giải cần
nắm rõ các tính chất khác nhau của mạng luồng, chẳng hạn đỉnh có dung lượng,
dung lượng có cận dưới và cận trên, đa cung, v.v. Đồng thời cần tổng kết kinh
nghiệm và phát huy tính sáng tạo.

\section{Phân tích ví dụ}

\subsection{Bài toán 1: quy hoạch phát triển dự án}

Công ty Macrosoft chuẩn bị lập kế hoạch phát triển trong tương lai. Các dự án
phát triển do các phòng ban đề xuất được tổng hợp thành một bảng kế hoạch.
Với mỗi dự án, bảng này cho biết khoản đầu tư cần thiết hoặc lợi nhuận dự
kiến. Vì việc thực hiện một số dự án phải phụ thuộc vào kết quả phát triển của
các dự án khác, nếu muốn thực hiện một dự án thì các dự án mà nó phụ thuộc
vào cũng không thể thiếu.

Bạn là chủ tịch công ty Macrosoft. Hãy chọn một số dự án sao cho công ty thu
được lợi nhuận ròng lớn nhất.

\paragraph{Dữ liệu vào.}
Dữ liệu gồm số dự án \(N\). Với mỗi dự án \(i\), có ngân sách \(C_i\) và tập
dự án phụ thuộc \(P_i\). Định dạng:
\begin{itemize}
  \item dòng đầu là \(N\);
  \item dòng thứ \(i\) tiếp theo mô tả dự án \(i\). Số đầu tiên là \(C_i\);
  số dương biểu thị lợi nhuận, số âm biểu thị đầu tư. Các số còn lại là chỉ
  số những dự án mà dự án \(i\) phụ thuộc.
\end{itemize}
Hai số kề nhau trên cùng dòng cách nhau bởi một hoặc nhiều dấu cách.

\paragraph{Dữ liệu ra.}
Dòng đầu là lợi nhuận ròng lớn nhất của công ty. Sau đó là một phương án chọn
dự án đạt lợi nhuận đó. Nếu có nhiều phương án, in phương án chọn ít dự án
nhất. Mỗi dòng in một chỉ số dự án, theo thứ tự tăng dần.

\paragraph{Giới hạn.}
\[
  0\le N\le 1000,\qquad -1000000\le C_i\le 1000000.
\]

\paragraph{Ví dụ.}
PDF gốc cho ví dụ gồm \(6\) dự án. Một nghiệm tối ưu có lợi nhuận ròng bằng
\(3\) và chọn các dự án \(1,2,3,4,6\).

\subsection{Trừu tượng hóa bài toán}

Cho đồ thị có hướng \(G=(V,E)\) gồm \(N\) đỉnh. Mỗi đỉnh đại diện cho một dự
án và có trọng số \(C_i\), biểu thị ngân sách của dự án. Dùng cung có hướng để
biểu diễn quan hệ phụ thuộc: cung từ \(u\) đến \(v\) nghĩa là dự án \(u\) phụ
thuộc vào dự án \(v\).

Bài toán trở thành: tìm một tập con \(V'\) của các đỉnh sao cho với mọi cung
\(\langle u,v\rangle\in E\), nếu \(u\in V'\) thì \(v\in V'\), và tổng trọng số
các đỉnh trong \(V'\) là lớn nhất.

\subsection{Tìm kiếm}

Nếu liệt kê mọi tập con của \(V\) thỏa điều kiện, độ phức tạp là \(O(2^n)\),
tức cấp số mũ. Dù cắt tỉa thế nào, vẫn khó thoát khỏi bản chất không đa thức.

\subsection{Quy hoạch động}

Cấu trúc của bài là đồ thị có hướng không chu trình chứ không phải cấu trúc
cây, nên không phù hợp với quy hoạch động. Nếu cố làm quy hoạch động, thực
chất gần giống tìm kiếm; vì số trạng thái rất lớn, độ phức tạp thời gian vẫn
là cấp số mũ.

\section{Mô hình luồng mạng}

\subsection{Cách xây dựng mạng}

Tạo \(N\) đỉnh tương ứng với \(N\) dự án, rồi thêm nguồn \(s\) và đích \(t\).
Nếu dự án \(i\) phải phụ thuộc vào dự án \(j\), thêm một cung từ \(i\) đến
\(j\) có dung lượng vô cùng lớn.

Với mỗi dự án \(i\):
\begin{itemize}
  \item nếu ngân sách \(C_i\) dương, tức dự án sinh lợi, thêm cung từ nguồn
  \(s\) đến đỉnh \(i\) có dung lượng \(C_i\);
  \item nếu \(C_i\) âm, tức dự án cần đầu tư, thêm cung từ đỉnh \(i\) đến đích
  \(t\) có dung lượng \(-C_i\).
\end{itemize}

Tìm lát cắt nhỏ nhất \((S,T)\) của mạng này, giả sử dung lượng là
\[
  C(S,T)=F.
\]
Gọi \(R\) là tổng ngân sách của mọi dự án sinh lợi, tức cận trên của lợi
nhuận ròng. Khi đó \(R-F\) là lợi nhuận ròng lớn nhất. Các đỉnh nằm trong
\(S\) biểu thị những dự án được chọn trong phương án tối ưu.

Trong ví dụ của PDF gốc, lát cắt nhỏ nhất có
\[
  S=\{s,1,2,3,4,6\},\qquad T=\{5,t\},
\]
và \(C(S,T)=5\). Vì \(R=8\), lợi nhuận ròng là
\[
  R-C(S,T)=8-5=3.
\]

\section{Chứng minh tính đúng đắn}

\subsection{Tương ứng giữa phương án và lát cắt}

Mỗi phương án chọn dự án có thể tương ứng với một lát cắt \((S,T)\) trong
mạng:
\[
  S=\{s\}\cup\{\text{các dự án được chọn}\},\qquad T=V-S.
\]

Với một phương án không thỏa quan hệ phụ thuộc, lát cắt tương ứng có đặc điểm:
tồn tại một cung dung lượng vô cùng lớn \(\langle u,v\rangle\), trong đó
\(u\in S\) nhưng \(v\in T\). Khi đó dung lượng lát cắt là vô cùng lớn, rõ ràng
không thể là lát cắt nhỏ nhất của mạng.

\subsection{Dung lượng lát cắt và lợi nhuận}

Xét một lát cắt bất kỳ tương ứng với một phương án hợp lệ. Lợi nhuận ròng của
phương án đó bằng
\[
  R-C(S,T).
\]
Những nguyên nhân khiến lợi nhuận thực tế nhỏ hơn cận trên \(R\) là:
\begin{enumerate}
  \item Không chọn dự án sinh lợi \(i\), tức đỉnh \(i\) nằm trong \(T\). Khi
  đó cung từ nguồn \(s\) đến \(i\) có dung lượng \(C_i\) bị cắt.
  \item Chọn dự án đầu tư \(i\), tức đỉnh \(i\) nằm trong \(S\). Khi đó cung
  từ \(i\) đến đích \(t\) có dung lượng \(-C_i\) bị cắt.
\end{enumerate}

Dung lượng \(C(S,T)\) chính là tổng dung lượng của hai loại cung trên. Vì vậy
lát cắt càng nhỏ thì lợi nhuận ròng của phương án càng lớn.

\subsection{Cách tìm lát cắt nhỏ nhất}

Theo định lý max-flow min-cut, lát cắt nhỏ nhất của mạng có thể tìm bằng thuật
toán luồng cực đại.

\section{Điểm mấu chốt của mô hình}

Mấu chốt của bài này nằm ở việc xây dựng mô hình toán học của mạng luồng. Điểm
đặc sắc là: các bài luồng mạng trước đây thường dùng giá trị luồng để biểu
diễn phương án, còn bài này dùng lát cắt để biểu diễn phương án, rồi khai thác
đầy đủ tính chất của lát cắt. Luồng chỉ là phương tiện để tìm lát cắt nhỏ
nhất.

Cách nhìn này mở ra một hướng mới để xây dựng mô hình luồng mạng cho bài toán.
Nhìn lần đầu, việc liên hệ bài này với luồng mạng là khá khó. Người giải phải
nắm vững các tính chất của mạng luồng, rồi qua nhiều lần so sánh và thử mô
hình mới có thể phát hiện điểm chung giữa chúng.

\section{Liên tưởng}

Một biến thể của bài toán là: ước lượng thêm thời gian hoàn thành cho mỗi dự
án và giả sử công ty chỉ có thể thực hiện một dự án tại một thời điểm. Khi đó
câu hỏi trở thành: làm thế nào để chọn một số dự án có thể hoàn thành trong
thời gian cho trước, sao cho công ty thu lợi lớn nhất. Tác giả cho biết đến
thời điểm viết bài vẫn chưa tìm được thuật toán hiệu quả cho biến thể này và
hy vọng các bạn quan tâm cùng nghiên cứu.

\section{Kỹ thuật lập trình}

\paragraph{Cấu trúc dữ liệu.}
Nên dùng danh sách kề.

\paragraph{Biểu diễn trực tiếp bài toán gốc.}
Ưu điểm là tiết kiệm bộ nhớ. Khuyết điểm là độ phức tạp lập trình lớn và tính
tổng quát kém.

\end{document}
